\documentclass[10pt, a4paper]{scrartcl}
%\usepackage{a4wide}					%Platzintensiver
\usepackage[german,ngerman]{babel}	%Deutsche Silbentrennung
\usepackage[utf8]{inputenc}		%Damit kann man auch Umlaute direkt eintippen
\usepackage{textcomp}
\usepackage{amsmath}
\usepackage{amssymb}
\usepackage{amsbsy}
\usepackage{amsfonts}
\usepackage{amsthm}
\usepackage{graphicx}	%including images
\usepackage{fancyhdr}	%for editiable headers and footers
\usepackage{setspace} 	%for 1,5times row space
\usepackage[colorlinks,pdfpagelabels,pdfstartview = FitH,bookmarksopen = true,bookmarksnumbered = true,linkcolor = black,plainpages = false,hypertexnames = false,citecolor = black] {hyperref}	
			%clickable TOC
\usepackage{multirow}	%linebreak in Table
\usepackage{enumerate}
\usepackage{float}
\usepackage{colortbl}	%colorTables
\usepackage{framed}   %begin{framed}
\usepackage{longtable} %begin{longtable}
\usepackage{makeidx}
\usepackage{natbib}
\usepackage{epigraph}
%\usepackage{listings}
\lstset{
language=php,                		% choose the language of the code
basicstyle=\footnotesize,       % the size of the fonts that are used for the code
numbers=left,                   % where to put the line-numbers
numberstyle=\footnotesize,      % the size of the fonts that are used for the line-numbers
stepnumber=1,                   % the step between two line-numbers. If it's 1 each line will be numbered
numbersep=5pt,                  % how far the line-numbers are from the code
backgroundcolor=\color{white},
showspaces=false,
showstringspaces=false,
showtabs=false,
frame=single,
tabsize=4,
captionpos=t,
breaklines=true,
breakatwhitespace=false,
%escapeinside={\%*}{*\%}{//},  	% if you want to add a comment within your code
keywordstyle=\color{green}\bfseries,
commentstyle=\color{blue},
morecomment=[l]{//},
morecomment=[s]{/*}{*/},
morestring=[b]",
morestring=[d]’,
frame=tb
}
\lstset{index={}}

%	Eigenes PHP Syntax-HL
\usepackage{listings}
\usepackage[dvipsnames]{xcolor}
\lstloadlanguages{Ruby}
\lstset{%
basicstyle=\ttfamily\color{black},
commentstyle = \ttfamily\color{gray},
keywordstyle=\ttfamily\color{blue},
stringstyle=\color{red},
showspaces=false,               % show spaces adding particular underscores
showstringspaces=false, 
frame=single
}



\newcommand{\ownTitle}{Diplomarbeit}
\newcommand{\ownTitleZ}{Testgetriebene Entwicklung von Web-Serveranwendungen auf der Basis von Ruby on Rails}
\newcommand{\ownAutor}{Stefan Wienert, HTW Dresden}



%%%%%%%%%%%%%%%%%%%%%%%%%%%%%%%%%%
%% HEAD AND FOOTLINES 
%%%%%%%%%%%%%%%%%%%%%%%%%%%%%%%%%%
\pagestyle{fancy}
\renewcommand{\sectionmark}[1]{\markboth{\thesection\ #1}{}}
\renewcommand{\subsectionmark}[1]{\markright{\thesubsection\ #1}}
\lhead[ \leftmark   ]{\textbf{\ownTitle}}
\rhead[\textbf{\ownTitleZ}]{\leftmark}
\lfoot[\thepage    ]{\scriptsize \textcopyright2011, \ownAutor}
\cfoot[]{}
\rfoot[\scriptsize \textcopyright2011, \ownAutor]{\thepage}

%%%%%%%%%%%%%%%%%%%%%%%%%%%%%%%%%%%%%%%
%  Macros
%%%%%%%%%%%%%%%%%%%%%%%%%%%%%%%%%%%%%%%%
\renewcommand{\footrulewidth}{0.5pt}
\newcommand{\cbox}{\includegraphics[width=0.2cm]{material/cbox.png} \ }
\definecolor{Gray}{rgb}{0.85,0.85,0.90}
% \setcounter{tocdepth}{2}	%kleineres TOC
\newcommand{\arr}{$ \Longrightarrow$}

\makeindex 	%start Aufzeichnun eines Stichwortverzeichnisses

%%%%%%%%%%%%%%%%%%%%%%%%%%%%%%%5
%% Metainfos
%%%%%%%%%%%%%%%%%%%%%%%%%%%%%%%


\title{\ownTitle}
\author{\ownAutor}
\pdfinfo{
 /Title  (\ownTitle)
 /Author (\ownAutor)
 /Subject (\ownTitleZ)
 /Keywords()
}
\date{\today}
\onehalfspacing

\setlength{\epigraphwidth}{.8\textwidth}
%%%%%%%%%%%%%%%%%%%%%%%%%%%%%%%%%%%%%%%%%%
\begin{document}
 %\maketitle
 %\thispagestyle{empty}
\begin{center}
\begin{tabular}{lcr}
 \Large{HTW Dresden} & \verb|       |& \multirow{3}{*}{\includegraphics[height=1.353cm]{material/htwlogo.jpg}} \\
 \Large{Fachbereich Informatik} &  & \\
\ownTitle &  & \\
\end{tabular}\end{center}
\begin{center}


\end{center}
\begin{verbatim}





\end{verbatim}
\begin{center}
\textbf{\LARGE{\ownTitle}}

"`\ownTitleZ"'

\end{center}
\begin{verbatim}



\end{verbatim}
\begin{center}
%\textbf{\LARGE{Wahrscheinlichkeitsrechnung und Statistik}}
\end{center}
\begin{verbatim}








\end{verbatim}
\begin{flushleft}
\begin{tabular}{lll}
\textbf{Autor:} & & Wienert, Stefan\\
& & \\
\textbf{Seminargruppe} & & 07/041/02\\
& & \\
& & \\
\textbf{Betreunder Professor:} & & Prof. Dr. Fritzsche\\
& & \\
& & \\
\textbf{Datum:} & & \today\\

\end{tabular}

\end{flushleft}
 %\tableofcontents		%Inhaltsverzeichnis
 \section{Motivation}


Kurz nach der Firmengründung der pludoni GmbH absolvierte der Autor dort sein Pflichtpraktikum, und war bis zum heutigen Tag als Werkstudent tätig.
Währenddessen nahmen Programmierer verschiedener Erfahrungsstufen und insbesondere auch Praktikanten an der Neu- und Weiterentwicklung der Software teil. Dies hat zur Folge, dass die Komplexität der Software inzwischen ein Level erreicht hat, dass das Maß an Regressionsfehlern\footnote{fehlerauslösender Quelltextänderungen} steigt. Eine weitere wahrscheinliche Ursache dieses Problems ist, neben der fehlenden Erfahrung der Praktikanten, das Fehlen von automatisierten Tests.
Für ein neues Projekt, it-jobs-und-stellen.de, soll dies nun mit einem anderen Ansatz verlaufen. 

Neben der Umstellung auf ein modernes Web-Framework, sollen nun Tests im Einklang zum Code erstellt werden, um auf Knopfdruck eine umfassende Information über den Systemzustand zu erhalten, wie sie ein manueller Test in der Gründlichkeit niemals erreichen kann.

\subsection{Die pludoni GmbH}

Die pludoni GmbH ist ein junges dynamisches Dresdner Unternehmen, dass sich zum Ziel gesetzt hat, lokale Job-Communitys aufzubauen und zu betreuen, sowie Tools für die Vereinfachung der Personalarbeit mittelständischer Unternehmen zu entwickeln. Einige Beispiele für diese Communitys sind zur Zeit ITsax.de, ITmitte.de und MINTsax.de\footnote{http://www.itsax.de/, http://www.itmitte.de, http://www.mintsax.de/}.
\marginline{\includegraphics[width=0.8\marginparwidth]{material/pludoni.png}\\pludoni GmbH}
\paragraph{Funktionsweise der Jobcommunitys}
Die Jobcommunitys bestehen jeweils aus einer Anzahl meist mittelständischer Unternehmen einer Branche. Für ITsax.de ist das die IT-Branche. Neben diesem Branchenfokus sammeln sich auch nur Unternehmen einer spezifischen Region. Bei ITsax.de ist dies Mittel- und Ostsachsen, bei ITmitte.de z.B. Mitteldeutschland, d.h. Thüringen, Sachsen-Anhalt und Westsachsen. 
Diese Unternehmen, die einen jährlichen Mitgliedsbeitrag an die pludoni GmbH zahlen, dürfen Ihre für die Region relevanten Jobangebote auf dem jeweiligen Portal einstellen. Was die Jobcommunitys von pludoni von der der Konkurrenz unterscheidet, ist das sogenannte \textbf{Empfehlungssystem}. 
%TODO Bild
Viele der Personalchefs der beteiligten Firmen haben dieselbe Erfahrung gemacht, dass sie sehr guten Bewerbern absagen mussten, weil z.B. die Stelle schon vergeben wurde, die Fähigkeiten des Bewerbers nicht den Bedürfnissen des Unternehmens entsprachen, eine Einstellung verhinderte. Hier setzt pludoni mit seinen Jobcommunitys an, und stellt eine Infrastrukur zur gegenseiten Empfehlung dieser guten Bewerber bereit.  Ausgezeichnete Bewerber erhalten neben der Absage einen Empfehlungscode, mit dem sie sich auf dem Online-Jobportal bei einer der anderen Mitgliedsfirmen bewerben können. Die Software löst intern den Empfehlungscode auf und bestätigt dieser Firma nun, dass der Bewerber empfohlen wurde. 

Dieses Empfehlungssystem überzeugt die beteiligten Unternehmen. Aktuell wurden im letzten Jahr z.B. auf ITmitte.de über 800 Bewerberungen über das Portal versendet, von denen mehr als die Hälfte (440) mit Empfehlungscodes versehen waren \citep{joerg_klukas_startseite_2011}. Dies motivierte mittlerweile 54 Firmen bei den drei pludoni Communitys teilzunehmen \citep{joerg_klukas_referenzen_2011}.

\subsection{Arbeitsablauf in der pludoni GmbH}

Da der pludoni GmbH gegenwärtig weniger Büroplätze zu Verfügung stehen, als sie Mitarbeiter hat, findet ein Großteil der Arbeit in Heim- oder Telearbeit statt.
Zur persönlichen Abstimmung findet aber mindestens einmal pro Woche ein Meeting statt, in welcher sich 2-4 der Mitarbeiter treffen, um alte Aufgaben abzunehmen und neue zu diskutieren. Die Abnahme erfolgt dabei durch den Chef Jörg Klukas.

Zentrales Kommunikationsmittel der pludoni GmbH ist, neben der E-Mail, die Online-Aufgaben- und Fehlerverwaltung, Redmine\footnote{http://www.redmine.org - ein webbasiertes Projektmanagement-Tool auf der Basis von Ruby on Rails. Redmine kann für Benutzer- und Projektverwaltung, Diskussionsforen, Wikis, zur Ticketverwaltung oder Dokumentenablage genutzt werden, \textit{Wikipedia}}. Dort werden alle Aufgaben und Fehler erfasst und an die zuständigen Personen verteilt. 
Neben den technischen Aufgaben der Entwickler, werden auch nicht-technische Aufgaben der anderen Mitarbeiter verwaltet, wie z.B. die Gewinnung neuer Partner (Akquise) oder administrative Aufgaben. Für Redmine sprach, dass es:
\begin{itemize}
\item  eine einfache Bedienung,
\item  einen hohen Funktionsumfang erweiterbar durch Plugins, wie z.B. auch Zeitabrechnung und Projektplanung,
\item  ein E-Mail-Interface, welches über neue Tickets informiert, aber auch ein Antworten auf Tickets erlaubt,
\item  eine Integration von SCM\footnote{Source-Code-Management. Versionsverwaltung}, Subversion und GIT \end{itemize}
hat.

Trotz dieses Tools ist eine Heimarbeit aber immer mit Nachteilen in der Kommunikation behaftet. Dies hat bei den Programmierern teils gravierende Auswirkungen auf die Produktivität. Einerseits, da mitunter Code-Teile anderer abwesender Programmierer benutzt werden müssen,  andererseits, dass gleiche Funktionalität doppelt implementiert wird, weil der Überblick fehlt. Auch deswegen ist es leicht, bei Arbeiten am Code neue Fehler einzuführen. 
Für das kommende Projekt soll nun eine großflächige Testinfrastruktur erstellt werden, um das Regressionsrisiko zu minimieren und durch die Tests eine gute Quelltextdokumentation zu erhalten.


\subsection{Begriffsdefinitionen}
\begin{description}
 \item[Test] ist eine, meist automatisierte, Prüfung des Programmverhaltens bei definierten Eingabeparametern
 
 \item[Code-Qualität] beinhaltet die Qualitäten Lesbarkeit, Testbarkeit, Wartbarkeit, Erweiterbarkeit, Geringe Komplexität
 \item[Metrik] Eine Softwaremetrik ist eine statische oder dynamische Codeanalyse zur Generierung von statistischen Informationen über den Source-Code. Beispiele: Testabdeckung, Anzahl Codezeilen, Anzahl Bad Smells pro Codezeile.
 
 \item[Testabdeckung] auch: Testfallabdeckung oder Code-Coverage. Eine dynamische Code-Metrik die angibt, welche Codezeilen durch keinen Test abgedeckt wurde. Es wird unterschieden in die Stufen C0, C1 und C2 mit steigender Komplexität der Messung.\\
 C0: Messung jeder Zeile, ob diese ausgeführt wurde\\
 C1: Messung jedes Zweigs jeder Zeile, ob dieser ausgeführt wurde\\
 C2: Messung jedes möglichen Codepfades, ob dieser ausgeführt wurde
 %TODO Quelle http://blog.abakas.com/2008/04/code-coverage-complexity.html
 \item[Bad Smell] oder Code-Smell. Ist ein Anzeichen für eine suboptimale Code-Stelle, die auch ein Hinweis auf ein größeres Design-Problem sein kann. Oft auch ein Kandidat für ein Refaktoring
 \item[TDD] Test-Driven-Development, bezeichnet den Prozess der Testgetriebenen Entwicklung
 \item[BDD] Behavior-Driven-Development (Verhaltensgetriebene Entwicklung). Prinzipiell eine Umformulierung von TDD zur Unterstützung der Businessprozesse 
 %\item[ATDD] Acceptance-Test-driven-development. Im Gegensatz zu Unittests beim reinen TDD, stehen die Akzeptanztests
\end{description}



 \part{Theoretische Grundlagen}
\label{sec:theory}

\section{Die Programmiersprache Ruby}

Ruby ist eine Programmiersprache, die ab 1993 von Yukihiro Matsumoto entwickelt wurde. Dabei ließ er sich von seinen Lieblingsprogrammiersprachen Perl, Smalltalk, Eiffel, Ada und Lisp inspirieren, um eine neue Programmiersprache zu entwickeln, die sowohl funktionale und imperative Programmierung ermöglicht \citep{ruby_visual_identity_team_about_2011}. 

Eine vollständige Einführung in Ruby zu geben, würde den Rahmen dieser Diplomarbeit sprengen, weswegen ich mich auf die Herausstellung der Hauptmerkmale und Unterschiede zu anderen Sprachen konzentrieren werde, und was die Auswirkungen auf das Testgeschehen sind.

%Dem geneigten Leser seien zum Selbststudium 
% http://tryruby.org/
% .... TODO

\subsection{Einführung in Ruby}
\marginline{\includegraphics[width=0.8\marginparwidth]{material/ruby.png}
Ruby ist eine Multiparadigma-Sprache}
Ruby ist eine Multiparadigma-Sprache, die Objektorientierung, prozedurale, funktionale und nebenläufige Programmierung unterstützt. Im Gegensatz zu Java ist Ruby wie Smalltalk vollständig objektorientiert. So sind auch alle Datentypen wiederrum Objekte, was auch die primitiven Datentypen wie String und Integer umfasst. 
% TODO Quelle

Ruby ist dynamisch stark typisiert, d.h. dass die Zuweisung des Typs einer Variable zur Laufzeit des Programms geschieht. Der Typ einer Variable ergibt sich damit aus ihrem Wert.

Ruby ist eine interpretierte Sprache, auch Skriptsprache genannt. Dies heisst, dass der Programmcode zur Laufzeit analysiert und ausgeführt wird. 
% TODO Quelle!!

Ruby nimmt für sich in Anspruch, eine Sprache für Menschen, und nicht für Maschinen zu sein.

\setlength{\epigraphwidth}{\marginparwidth}
\marginline{\epigraph{Ruby is simple in appearance, but is very complex inside, just like our human body}{Yukihiro Matsumoto}}
\setlength{\epigraphwidth}{0.8\textwidth}

Im Nachfolgenden einige Beispiele für die Verwendung von Ruby, insbesondere die "`Alles ist ein Objekt"'-Philosophie.

\begin{lstlisting}[language=Ruby,label=Ruby Beispiele,caption=Ruby Beispiele]
 
>> puts "Hello World"
=> Hello World

>> 2.even?
=> true

>> "hallo".upcase
=> "HALLO"

>> Date.today + 2
=> #<Date: 2011-06-30>

>> a = 4 + Math.sqrt(9)
=> 7.0

>> if (0..10).include? a
>>   puts "a liegt zwischen 0 und 10"
>> end
=> a liegt zwischen 0 und 10

\end{lstlisting}

Desweiteren sei die sehr gute Unterstützung bei der bearbeitung von Strings hervorzuheben. Insbesondere die Regulären Ausdrücke sind denen von Perl beinahe gleichmächtig und fest in die Sprache eingebaut.

\begin{lstlisting}
>> match = "2011-06-20".match /(?<year>\d{4})-\d{2}-\d{2}/
>> puts match[:year]
=> 2011
\end{lstlisting}

Auch die Bearbeitung von Arrays und listenähnlichen Strukturen ist sehr bequem, dank der Verwendungsmöglichkeit von anonymen Funktionen, bei Ruby "`Blöcke"' genannt.

\begin{lstlisting}
>> [5,5,7,3].sort
=> [3, 5, 5, 7]

# Es kann auch eine benutzerdefinierte Sortierfunktion
# angegeben werden
>> [ "string",  "rails",  "ruby" ].sort_by{|item| item.length }
=> ["ruby", "rails", "string"]

# Die Quadratzahlen von 1 bis 5
>> (1..5).map{|element| element * 2}
=> [2, 4, 6, 8, 10]
\end{lstlisting}
Neben einem soliden objektorientierten System, bietet Ruby viele funktionale Aspekte um die Arbeit mit Arrays und Strings sehr einfach zu gestalten. Wichtig sei auch noch die Fähigkeit der Metaprogrammierung anzumerken.

Im Gegensatz zu Java oder C\# dürfen Klassen zur Laufzeit um Funktionen erweitert, oder alte sogar überschrieben werden. So ist es z.B. möglich, die String-Klasse um eigene Funktionen zu erweitern.

\begin{lstlisting}
>> class String
>>   def remove_whitespace
>>     self.gsub(/\s+/, "")
>>   end
>> end

>> "Dies ist ein Test".remove_whitespace
=> "DiesisteinTest"

\end{lstlisting}

Diese Beispiele sollten als kurzer Einstieg in Ruby dienen, und einen Querschnitt durch die Besonderheiten der Sprache aufzeigen.

Für eine weiter Vertiefung sei das Buch "`Programming Ruby 1.9"' empfohlen, das im Detail auf die neuste Version der Programmiersprache eingeht.

\subsection{Diskussion}

\paragraph{Ruby als Skriptsprache}
Der Vorteil, der Übersetzung zur Laufzeit, ist eine hohe Plattformunabhängigkeit und ein leichterer Buildprozess, da das Kompilieren entfällt. Allerdings hat dies in der Regel gravierende Geschwindigkeitsnachteile. 
Verfechter dynamischer Sprachen erklären weiterhin, dass diese ideal für Prototypische Implementierungen geeignet sind, da sich Anforderungen ständig ändern können. Weiterhin hätten Programme dynamischer Sprache eine potenziell hohe Wiederverwendbarkeit und eine höhere Lesbarkeit. 
%TODO Benchmark, Diskussion zur Architektur
%TODO Gegenbeispiele Groupon, Twitter

Desweiteren bleiben Fehler, die der Compiler bereits entdeckt hätte, bis zur Ausführung oder schlimmstenfalls noch länger unentdeckt. Dazu gehören z.B. Tippfehler, bei denen der Wert einer nicht deklarierten Variable ausgelesen wird. Im Gegensatz zu z.B. PHP, wirft Ruby aber dann eine Exception.

Auf das Testen hat dies eine direkte Auswirkung. Viele Quellen belegen, dass eine dynamisch typisierte Sprache mehr Tests benötigt, als eine statisch typisierte. 
%TODO Quelle http://stackoverflow.com/questions/125367/dynamic-type-languages-versus-static-type-languages



\paragraph{Metaprogrammierung}
Die Fähigkeiten zur Metaprogrammierung von Ruby bieten vielerei Möglichkeiten um Probleme effektiv zu lösen, die andernfalls nur mit erheblichem Aufwand, oder gar nicht zu lösen sind. So verwendet das beliebte Objektrelationale Datenbankframework ActiveRecord Metaprogrammierung, um einfache SQL-Statements zu erstellen.
\begin{lstlisting}
>> Person.find_by_first_name("Stefan")
  Person Load (0.2ms)  SELECT persons.* FROM persons
    WHERE users.first_name = 'Stefan' LIMIT 1
\end{lstlisting}

Die Methode \texttt{find\_by\_first\_name} existiert nicht, und wird zur Laufzeit auf Basis des Namens gebaut.

Auch das sehr beliebte Testframework Rspec, auf das wir später kurz eingehen werden, verwendet Metaprogrammierung, um Testfälle und Zusicherungen wie fast in der englischen Sprache zu formulieren.

\begin{lstlisting}
>> 4.should == 3
RSpec::Expectations::ExpectationNotMetError: expected: 3
\end{lstlisting}

All diese Methoden können, richtig angewendet, zur Verbesserung der Lesbarkeit der Programme, und damit zur Erhöhrung der Wartbarkeit, führen.
Falsch angewendet stellen sie jedoch eine Gefahr da. Beim Verwenden von externen Bibliotheken, oder in einem großen Entwicklerteam können dadurch kuriose Fehler auftreten die nur sehr schwer zu finden sind. Potenziell jedes Quelltextfile kann jede Klasse zur Laufzeit verändern, ohne dass es eine Warnung gibt. Umgehen kann man dieses Problem bestenfalls durch das konsequente Einordnen in Namensräume, bei Ruby Module genannt.
%TODO Quelle?

\paragraph{Schlussfolgerung}

Die Verwendung von Ruby und anderen dynamischen Sprachen birgt durchaus Risiken, die zu beachten sind. Falls man diese Risiken im Kopf behält, und die Möglichkeiten der Sprache nutzt, um die Lesbarkeit zu verbessern, sind sie gerechtfertigt. Gerade bei der Entwicklung kleinerer Entwicklerteams oder Projekten mit engem Budget können dynamische Sprachen ihre Vorteile ausspielen, da sie eine eine schnellere Entwicklung ermöglicht. Im Gegensatz zu den meisten auf C basierten Sprachen, ist die Syntax von Ruby äußerst leserlich, da nur wenige Sonderzeichen verwendet werden. Auch biete Ruby mehr Funktionalität pro Programmzeile, da die Deklaration entfällt und es viel sogenannten syntaktischen Zucker gibt. Auch dies kann, richtig angewendet, der Lesbarkeit zuträglich sein.
% TODO QUelle hier was?
\epigraph{Sometimes people jot down pseudo-code on paper. If that pseudo-code runs directly on their computers, it's best, isn't it? Ruby tries to be like that, like pseudo-code that runs. }{Yukihiro Matsumoto}

\subsection{Ruby on Rails}

Für das Projekt IT-jobs-und-stellen.de soll das Webframework Ruby-on-Rails verwendet werden. 
Dieses erlangte seit 2006 einer wachsenden Beliebtheit, was sich z.B. im Stellenmarkt widerspiegelt.
%TODO Quelle http://www.businessinsider.com/heres-why-ruby-on-rails-is-hot-2011-5

Viele professionelle Websites, die meist als Startup begannen, setzen bis heute auf Rails. Darunter z.B. Yellow Pages, die Gelben Seiten der USA, Github, eine sehr beliebte Community für OpenSource Programmierer,  Groupon, dem führenden Unternehmen bei Online-Gutscheinen und XING, einer deutschen Online-Community für Business-Kontakte.
% TODO http://rubyonrails.org/applications

Im Folgenden werden die Grundzüge von Ruby on Rails näher erläutert.


\section{Softwaretests}

\subsection{Warum testen}

(Softwarefehler) \url{http://www.nfranze.de/download/Diplomarbeit_Nico_Franze.pdf}
  Zusammenfassende Begriffsdefinition

\subsection{Arten von Tests}

\subsection{Testtools in Ruby}
\subsubsection{Test::Unit und Minitest}
Test::Unit (Ruby 1.8.7) und Minitest (1.9.2) sind die Testbibliotheken, die Ruby standardmäßig mitbringt. Beide basieren auf dem xUnit, bzw. SUnit Design von Kent Beck, und sind für Nutzer von JUnit oder NUnit leicht nachvollziehbar.

Für eine zu testende Klasse wird eine analoge Testklasse erstellt. Diese trägt per Definition denselben Namen wie die zu testende Klasse mit einem "`test"' am Anfang. Um z.B. eine Klasse "`job"' zu testen, wird eine Datei \texttt{test\_job.rb} (Ruby Standard) oder \texttt{job\_test.rb} (Rails Standard) erstellt. Dort wiederrum wird eine Klasse mit Namen \texttt{TestJob} definiert. 

Ein Beispieltest sieht z.B. so aus:
\begin{lstlisting}
require "job"

class TestJob < Test::Unit::TestCase
  def setup
    @job = Job.create
  end
  
  def teardown
    Job.delete_all
  end
  
  def test_job_exists
    @job.title =  "Ruby on Rails Entwickler
    @job.add_location_to_title( "Dresden")
    
    assert_equal( "Ruby on Rails Entwickler in Dresden",  Job.first.title)
  end
end
\end{lstlisting}
Unsere Klasse TestJob erbt von der TestUnit Basisklasse. Sie beinhaltet die Methoden "`setup"' und "`teardown"', die jeweils vor, respektive nach jedem einzelnen Testfall aufgerufen werden.
In der Setup-Methode nehmen wir z.B. das Anlegen eines Jobs vor, in der Teardown Methode löschen wir alle Jobs in der Datenbank, um einen sauberen Test zu gewährleisten

Danach können nun beliebig viele Testmethoden folgen, deren Namen mit \texttt{test\_} beginnen müssen.
Jede Testmethode besteht in der Regel aus einer Initialisierung (kann in die setup-Methode ausgelagert werden), der Ausführung einer zu testenden Aktion und dem Prüfen der danach geltenden Eigenschaften mittels Assertions. Diese Zusicherungen sind Prädikate die oft Gleichheit oder Boolsche Rückgabewerte prüfen.

%TODO Weitere Prädikate erwähnen

Rails erledigt das Anlegen und Löschen von Testdaten selbständig. Diese werden als Fixtures bezeichnet und extern definiert. Alternativ ist der Einsatz sogenannter Factories möglich, um schnell Objekte mit bestimmten Eigenschaften zu erstellen. In jedem Fall setzt Rails die Datenbank nach jedem einzelnen Test zurück.


\subsubsection{Rspec}
\subsubsection{Cucumber}

\section{Testgetriebene Entwicklung}
Testgetriebene Entwicklung, im Englischen Test-Driven-Development (TDD) wurde erstmalig von Kent Beck 2003 im Detail erläutert. Zuvor war die Technik "`Test-First"' aber schon seit 1999 im Kontext von Extreme Programming (XP) bekannt.
%TODO Quelle: "Extreme Programming". Computerworld. Retrieved January 11, 2011.
%TODO Quelle: Beck, K. Test-Driven Development by Example, Addison Wesley, 2003
Damit wurde aus der Entwicklungs-Methode "`Test-First"' ein Software-Entwicklungsprozess. 

Im Folgenden werde ich diesen Prozess näher beleuchten und am Beispiel von Ruby on Rails typische Testwerkzeuge aufzeigen. \subsection{Motivation}
  Das Erstellen einer gut abdeckenden Test-Suite für ein jedes größeres Softwareprojekt ist eine wichtige Vorraussetzungen um interne Qualitäten, wie Wartbarkeit und Zuverlässigkeit zu aktivieren. TDD soll nicht dazu dienen, die Software zu verifizieren. Dies ist aber ein positiver Nebeneffekt. Das Hauptziel ist es, den Code in Einklang mit dem Test zu schreiben, so dass der Test den Code antreibt (Test drives thes code). Der messbare Effekt davon, ist ein gut-testbarer Code, welcher in der Regel auch ein gut-wartbarer und verständlicher Code. Bad-Smells, wie God-Methode und geringe Kohäsion, werden schon im Keim erstickt werden, da sie selbst nur äußerst schwer zu testen sind.
  
  
  %Psychologische...
  % Zuversicht
  TDD hat auch eine psychologische Motivation.
  In Abbildung \ref{} ist ein Einflussdiagramm zu sehen, welches vereinfacht Zusammenfänge darstellen soll. Ein direkter Pfeil zeigt eine direkte proportionale Abhängigkeit, ein Pfeil mit einem Kreis eine indirekt proportionale Abhängigkeit. In diesem Diagramm soll der Zusammenhang zwischen psychischen Druck auf den Programmier, Testen und Defekte visualisiert werden. Druck führt somit zu weniger Tests, weniger Tests zu mehr Fehlern und diese wiederrum zu höherem Druck. Dieser klassische Teufelskreis der Softwareentwicklung, der gerade zum Ende eines Entwicklungsprozesses auftritt, soll durch TDD effektiv verhindert werden. 
  % TODO Figure A.5 
  TDD fördert die Entwicklung in kleinen Schritten, und ermöglicht durch bestande Tests kleine "`Belohnungen"' für den Programmierer. Dadurch ist es leichter einen gewissen Arbeitsrhythmus zu erhalten, was stellenweise dem "`Flow"'\footnote{Schaffen-, Tätigkeitssrausch}  ähnelt oder diesen strukturiert ergänzen kann.
  % TODO http://www.agilecoachjournal.com/post/Test-Driven-Development.aspx
  
\subsection{Ablauf}
  Ziel ist es, vor der Implementation eines Codes, einen Unittest zu implementieren. Davon ausgehend soll der geringstmögliche Code implementiert werden, damit der Test besteht. Zum Schluss wird refaktorisiert, bei TDD auch als Designphase genutzt.
  
  %TODO Abbildung mit Beschreibung
  Im Ganzen sind das also die Phasen:
  \begin{enumerate}
   \item Schreibe einen Test von dem Feature, das als nächstes entwickelt werden möchte\\Oder von dem Test, der als nächstes auf der Liste steht
   \item Red: Führe alle Tests aus, um sicherzugehen, dass der Test fehlschlägt. Andernfalls ist der Test überflüssig.
   \item Green: Nachdem der Test fehlschlägt, implementiere nun den leichtesten Code, der den Test besteht\\
   Dies kann ausdrücklich auch eine Fake-Implementierung sein. Wichtig ist, dass diese Phase so schnell wie möglich verlassen wird.
   \item Refactor: Nachdem der Test bestanden wird, folgt nun die \textbf{wichtigste Phase}, die Refaktorisierungsphase.\\
   Da wir bereits einen Test haben, der unser gewünschtes Systemverhalten widerspiegelt, können wir gefahrlos refaktorisieren, d.h. im Einzelfall Duplikation eleminieren. In dieser Phase findet das Design des Codes statt. Man macht sich Gedanken, wie die vorhanden Klassen optimal refaktorisieren können, um Code-Smells zu eleminieren.
  \end{enumerate}
  
  Jeder Unittest soll prinzipiell nur eine Eigenschaft testen, die Entwicklung erfolgt also in kleinen Schritten. Dies hat direkte Auswirkungen auf die zu entwickelnden Objekte und Methoden, die ebenfalls übersichtlich werden sollen, und somit dem Funktionsbegriff, eine Methode für eine Aufgabe und erweitert eine Klasse für eine Aufgabe, gerecht werden.
  
  
  Das ganze lässt sich auch in den übergeordneten Prozess zum Entwickeln eines Features einordnen.
  %RSPEC Buch macht das vor, Noch mal schauen, evtl alles Käse was ich hier geschrieben habe
  \begin{enumerate}
   \item Schreibe einen Integrations/Akzeptanztest um das aktuelle Feature zu implementieren
   \item Nach einer möglichen Überlegungs und Analysezeit, kann nun mit dem Implementierung der Teilschritt nach TDD Ablauf von oben verfahren werden
   \item Nachdem die Teilschritte implementiert wurden, kann nun mit dem Erfüllen des Akzeptanztests fortgefahren werden
  \end{enumerate}
  
  Somit werden 2 Testebenen erstellt, die Akzeptanz- und die Unittests.

  Im Prinzip soll jeder Änderung der Programmlogik ein fehlgeschlagener Test vorrausgehen. Auch bei der Behebung von Defekten, also dem Bugfixing, soll zuerst ein Test geschrieben werden, der genau das Verhalten des Bugs widerspiegelt, und erst dann gefixt werden.
  
    
  \subsection{Mocks und Stubs}
  Beim Testen im Allgemeinen und bei TDD im Besonderen wird das Testen durch externe Abhängigkeiten erschwert. Dies können z.B. Klassen sein, die noch gar nicht implementiert wurden, externe Ressourcen (Netzwerkzugriffe, Versenden von Mails) oder externe Prozessen (Bezahlen in einem Onlineshop) sein. In diesen Situationen ist es angebracht, auf sogenannte Mocks und Stubs zurückzugreifen.
  
  Ein Stub ist eine nachahmende Funktion oder Objekt, welches die schwer zu isolierende Klasse während des Testfalls ersetzt. Im Beispiel ein Bezahlprozess einer Bestellung.
  \begin{lstlisting}
def test_report_failed_payment
  Payment.stubs(:pay).returns(false)
  
  bestellung = Bestellung.new()
  bestellung.commit()
  
  assert bestellung.errors.present?
end


  \end{lstlisting}
  Mit dem oben angegeben (Pseudo) Rubycode würde man z.B. mittels des Mock-Frameworks mocha ein Mockobjekt erzeugen, welches den (Fantasie-) Bezahlprozess nachahmt, und die Methode "`pay"' ersetzt, so dass sie immer "`false"' zurückgibt.
  So kann z.B. das Objekt Bestellung gefahrlos Bezahlungen auslösen. Der genaue Ablauf innerhalb des Bezahlprozesses ist die Bestellung unwichtig, lediglich der Rückgabewert, ob die Bezahlung erfolgreich war oder nicht.
  
  Als Ergänzung dazu gibt es Mocks. Ähnlich wie die Stubs ersetzen sie Methoden oder Objekte, um statt komplexer Operationen fixe Werte zurückzugeben. Zusätzlich dienen Mocks selbst als Testfall. Ein Mock wartet darauf, ob die Methode, wie sie definiert wurde, auch tatsächlich aufgerufen wurde.
  
  Hier z.B. ein Mock, um statt des heutigen Datums das von vor einer Woche zurückzugeben
  \begin{lstlisting}
def test_always_fail
  Date.mocks(:today).returns( 7.days.ago)
end
  \end{lstlisting}
  Der gezeigte Programmcode wird jedesmal fehlschlagen, da von einem Mock erwartet wird, dass er während des Tests genau einmal aufgerufen wird. Ist dies nicht der Fall, gilt der Test als nicht bestanden. Mocks fungieren somit als zusätzliche Möglichkeit Interna des Programmflusses zu testen.

  %RSPEC Mock basiert
  
  \subsection{Vorteile von TDD}
  Produktiver, 
    da weniger manuelle Tests
    weniger Debugger notwendig
  
  
  
  \subsection{Nachteile von TDD}
  %\subsection{Bewertung}
  \subsection{Varianten: BehaviorDD, DesignDrivenTesting}
\section{Code-Metriken}
        \subsection{Notwendigkeit von Code Metriken}
        \subsection{Überblick über Code-Metriken und Skalen}
        \subsection{Testabdeckung}


 \part{Praktische Erfahrungen und Ergebnisse}
\label{sec:praxis}
 \part{Zusammenfassung und Ausblick}
\label{sec:fazit}

 \appendix

\chapter{Eigenschaften erfolgreicher Tests}

Das Vorhandensein von zahlreichen Tests reicht nicht aus, um die Qualität der Tests zu beuteilen. Stattdessen existieren einige Qualitätskriterien, um Tests auf ihre Sinnhafigkeit und Wartbarkeit zu untersuchen:

\begin{description}
 \item[Unabhängigkeit (Independence) und Isolation] Ein Test \index{Test}ist unabhängig, falls er nicht durch andere Tests beeinflusst wird. Auch die Reihenfolge, in der die Tests ausgeführt werden, darf auf das Ergebnis keinen Einfluss üben. Siehe auch \citep{beck_test_2002}, \citep{palermo_guidelines_2006} sowie \citep[Karte 45]{langr_agile_2011}.
 \item[Wiederholbarkeit (Repeatability] Ein Test \index{Test}wird als wiederholbar bezeichnet, wenn er mehrmals hintereinander ausgeführt werden kann, und dabei jedes mal dasselbe Ergebnis liefert \citep[Karte 45]{langr_agile_2011} \citep{rappin_rails_2011}. Problematisch sind dabei z.B. Datum und Zeit, sowie Zufallsfunktionen
 \item[Klarheit (Clarity)] Ein Test \index{Test}ist klar, wenn sein Zweck sofort verständlich wird \citep{rappin_rails_2011}, \citep{palermo_guidelines_2006}. Damit wird einerseits die Lesbarkeit gemeint. Anderseits schließt dies auch ein, ob der Test genau eine Eigenschaft testet und nicht redundant zu anderen Tests ist. Dies hat zur Folge, dass die Tests wartbarer werden und als Code Dokumenation dienen können.
 \item[Präzise (Conciseness)] Ein Test \index{Test}ist präzise, wenn er so wenig Code und so wenige Objekte wie möglich benötigt, um sein Ziel zu erreichen \citep{palermo_guidelines_2006} \citep{rappin_rails_2011}. Eine Auswirkung ist, dass der Test schneller wird.
 \item[Robustheit (Robustness)] Ein Test \index{Test}ist robust, wenn es eine direkte Korrelation zum zu testenden Code gibt: Ist der Code korrekt, so ist der Test erfolgreich. Ist der Code falsch, so schlägt der Test fehl. Nicht-robuste Tests werden auch "`zerbrechlich"' (brittle) genannt. Dazu zählen auch sogenannte tautologische Tests, die immer erfolgreich Verlaufen, und keine Aussage über den zugrunde liegenden Programmcode geben
 \item[Schnelligkeit] Gute Tests sollten schnell sein, da die Test-Suite\index{Test!Test-Suite} ansonsten deutlich seltener ausgeführt würd \citep{langr_agile_2011} \citep{palermo_guidelines_2006} \citep{nagappan_realizing_2008}. Siehe auch Präzision.
 \item[Zeitliche Nähe] Gute Tests entstehen in zeitlicher Nähe zu dem zugrunde liegenden Programmstück \citep{langr_agile_2011}. Innerhalb von TDD\index{TDD} ist dieses Kritierium immer erfüllt.
 \end{description}

% Microsoft Vorgehensweise:
%
% Start TDD from the beginning of projects. Do not stop in the middle and claim it
% doesn’t work. Do not start TDD late in the project cycle when the design has already
% been decided and majority of the code has been written. TDD is best done
% incrementally and continuously.
% For a team new to TDD, introduce automated build test integration towards the second
% third of the development phase—not too early but not too late. If this is a “Greenfield”
% project, adding the automated build test towards the second third of the development
% schedule allows the team to adjust to and become familiar with TDD. Prior to the
% automated build test integration, each developer should run all the test cases on their
% own machine.
% Convince the development team to add new tests every time a problem is found, no
% matter when the problem is found. By doing so, the unit test suites improve during the
% development and test phases.
% Get the test team involved and knowledgeable about the TDD approach. The test team
% should not accept new development release if the unit tests are failing.
% Hold a thorough review of an initial unit test plan, setting an ambitious goal of having
% the highest possible (agreed upon) code coverage targets.
% Constantly running the unit tests cases in a daily automatic build (or continuous
% integration); tests run should become the heartbeat of the system as well as a means to
% track progress of the development. This also gives a level of confidence to the team
% when new features are added.

% Encourage fast unit test execution and efficient unit test design. Test \index{Test}execution speed is
% very important since when all the tests are integrated, the complete execution can
% become quite long for a reasonably-sized project and when using constant test
% executions. Tests results are important early and often; they provide feedback on the
% current state of the system. Further, the faster the execution of the tests the more likely
% developers themselves will run the tests without waiting for the automated build tests
% results. Such constant execution of tests by developers may also result in faster unit
% tests additions and fixes.

% Share unit tests. Developers’ sharing their unit tests, as an essential practice of TDD,
% helps identify integration issues early on.

% Track the project using measurements. Count the number of test cases, code coverage,
% bugs found and fixed, source code count, test code count, and trend across time, to
% identify problems and to determine if TDD is working for you.

% Check morale of the team at the beginning and end of the project. Conduct periodical
% and informal surveys to gauge developers’ opinions on the TDD process and on their
% willingness to apply it in the future.

%
% % Einige dieser Punkte können durch Metriken überprüft werden. Dazu mehr im Abschnitt \ref{sec:metrics}.
% %  Code Craft S 144

\chapter{Nutzung von Cucumber in Verbindung mit Selenium für Firefox und Guard ohne X-Server}
Für ein effektives Test-Setup wurden folgende Testtools benutzt
\begin{description}
 \item[Guard] Automatische Testausführung bei Änderungen der Dateien
 \item[Xvfb] Ist ein X-Server, der ohne Grafikausgabe arbeitet. Er eignet sich also für die Ausführung von GUI-Programmen, wie Mozilla Firefox auf Remote-Servern.
 \item[Selenium\index{Selenium}] Interface zur Fernsteuerung von Browsern
 \item[Spork] Erhöht die Testausführungsgeschwindigkeit, da Teile von Rails\index{Ruby-on-Rails} im Hintergrund gehalten werden, und zwischen den Tests nicht neu geladen werden müssen
\end{description}

\subsection*{Installation}
Es muss bereits installiert sein: firefox, xvfb, Rails > 3.0
\begin{enumerate}
 \item In das Gemfile folgendes eintragen:
 \begin{lstlisting}[caption=RAILS\_ROOT/Gemfile]
group :test, :cucumber do
  gem "capybara"
  gem 'cucumber'
  gem "cucumber-rails"
  gem 'database_cleaner'
  gem "launchy"
  gem "guard"
  gem 'guard-cucumber'
  gem 'guard-spork'
  gem "rb-inotify"  # Für Linux
  gem "spork", :git => "git://github.com/timcharper/spork.git"
end

 \end{lstlisting}
 Danach folgendes ausführen:

 \begin{verbatim}
bundle install
rails g cucumber:install
 \end{verbatim}


 \item Eine Datei Guardfile anlegen. Hier kommen alle Anweisungen zur Steuerung von guard hinein
 \begin{lstlisting}[caption=RAILS\_ROOT/Guardfile]
group "cucumber" do
  guard 'spork' do
    watch('config/application.rb')
    watch('config/environment.rb')
    watch('features/support/env.rb')
    watch(%r{^config/environments/.+\.rb$})
    watch(%r{^config/initializers/.+\.rb$})
    watch('spec/spec_helper.rb')
  end
  guard 'cucumber', :cli => '--no-profile --color --format pretty --strict RAILS_ENV=test' do
    watch(%r{^features/.+\.feature$})
    watch(%r{^features/support/.+$})                      { 'features' }
    watch(%r{^features/step_definitions/(.+)_steps\.rb$}) { |m| Dir[File.join("**/#{m[1]}.feature")][0] || 'features' }
  end
end
\end{lstlisting}
\item In der Datei \texttt{config/cucumber.yml} die Option \texttt{--drb} für "`default"' und "`wip"' setzen.


\item Die Datei \texttt{features/support/env.rb} anpassen, um sie mit Spork kompatibel zu machen
\begin{lstlisting}[caption=features/support/env.rb]
require 'cucumber/rails'
require "spork"
require 'test/unit/testresult'
Test::Unit.run = true

Spork.prefork do
  ENV["RAILS_ENV"] ||= "test"
  require File.expand_path(File.dirname(__FILE__) + '/../../config/environment')
  require 'cucumber/formatter/unicode'
  require 'cucumber/rails'
  require 'test/unit/testresult'
  ActionController::Base.allow_rescue = false
  begin
    DatabaseCleaner.strategy = :truncation
  rescue NameError
    raise "You need to add database_cleaner to your Gemfile (in the :test group) if you wish to use it."
  end
end

Spork.each_run do
  require 'cucumber/rails/world'
end

\end{lstlisting}

\end{enumerate}


\subsection*{Nutzung}

\begin{enumerate}
 \item Aktivierung von Xvfb auf Displayport 99
 \begin{verbatim}
  Xvfb :99 -ac
 \end{verbatim}

 \item Aktuelle Shell auf diesen Displayport einstellen
 \begin{verbatim}
  export DISPLAY=:99
 \end{verbatim}

 \item Guard starten
 \begin{verbatim}
  guard -g cucumber
 \end{verbatim}
\end{enumerate}

Nun werden automatisch alle Cucumber\index{Cucumber}-Features getestet. Falls Selenium\index{Selenium} verwendet wird, dann wird der Firefox im Hintergrund gestartet, ohne dass dies sichtbar wird.



\newpage
\listoffigures

\addcontentsline{toc}{chapter}{Quellcode-Listings}
\lstlistoflistings
%\listoflistings


%  \cite{*}
\bibliographystyle{alphadin}    % Verbos

%  \bibliographystyle{natbib}     %

\renewcommand{\indexname}{Stichwortverzeichnis}
%\addcontentsline{toc}{chapter}{Stichwortverzeichnis}
\printindex

\end{document}
 
Vorlagen:
\begin{longtable}{|l|l|}
\hline
\rowcolor{Gray}
Falsch&Richtig\\
\hline
Studenten&Studierende\\
\hline
Abschlussprogramm&Degree Seeking Students\\
\hline
Doppeldiplom&Doppelabschluss[programm]\\
\hline
HTWD & HTW Dresden\\
\caption{Deutsche Terminologie}
\label{tb:dt}
\end{longtable}
